%% USPSC-Cap3-Conclusao.tex
% ---
% Conclusão
% ---
\chapter{Conclusão}
% ---
Este trabalho investigou a otimização de kernels de inferência para a operação
de autoatenção em Vision Transformers, com foco na redução de acessos à memória
global da GPU. A partir de um kernel customizado em OpenAI Triton, foram
exploradas técnicas de fusão de operadores, tiling e \textit{softmax online},
mantendo o escopo restrito ao forward de atenção densa não causal em cenário de
inferência.

Os resultados confirmam a hipótese principal do trabalho: evitar a materialização
da matriz intermediária de atenção reduz substancialmente o tráfego em DRAM e
melhora o desempenho em relação à implementação explícita em PyTorch. Na
configuração avaliada, o kernel Triton atingiu speedup de até 6,16$\times$ sobre
o eager em $N=1024$, enquanto o SDPA do PyTorch atingiu 9,19$\times$. O profiling
de baixo nível reforçou essa leitura: no mesmo ponto, o eager movimentou cerca de
297,3 MB em DRAM, enquanto Triton e SDPA movimentaram aproximadamente 4,79 MB e
4,94 MB, respectivamente.

A comparação com o SDPA também delimita a contribuição do trabalho. Embora o
kernel customizado reduza o tráfego de memória para um patamar semelhante ao da
biblioteca otimizada, seu throughput permanece inferior. Isso indica que a
eliminação de escritas intermediárias é condição importante, mas não suficiente,
para aproximar a execução dos limites efetivos do hardware. Fatores como
ocupação, pressão de registradores, reutilização local dos blocos, escalonamento
de warps e escolha dos parâmetros de tiling continuam determinantes para o
desempenho final.

Assim, a contribuição central não é propor um kernel superior ao estado da arte,
mas apresentar uma implementação experimental e um protocolo de análise capaz de
explicar onde os ganhos ocorrem e quais gargalos permanecem. Como trabalhos
futuros, destacam-se a aplicação de autotuning aos parâmetros do kernel, a
avaliação em diferentes GPUs e tipos numéricos, a inclusão de atenção causal e a
extensão para treinamento, incluindo a etapa backward.
